\section{Proposed Approach}
\label{sec:proposed-approach}

The project will establish a TorchInductor baseline, profile GNN workloads across graphs with different structures, and use the results to choose operations for specialization. The integration mechanism will be chosen after profiling.

\subsection{Establishing Baselines}

The first stage will establish reproducible inference workloads using PyTorch Geometric. Each workload will specify its model configuration, input graph, feature dimensions, and numerical data type. Model parameters will remain fixed during inference.

The primary baseline will be PyTorch Geometric workloads compiled with \texttt{torch.compile} using TorchInductor. Uncompiled PyTorch Geometric execution will also be measured for context. Compilation failures or graph breaks, where execution leaves the captured computation and returns to ordinary execution, will be recorded so that comparisons describe what actually ran.

The initial set will contain a few representative GNN models that use neighborhood aggregation and dense feature transformations. The models will be chosen based on compatibility with the available software and hardware.

\subsection{Investigating Topology Effects}

The project will profile these workloads across graphs with different structural characteristics. Synthetic graphs will allow selected properties to vary under controlled conditions. Real-world graphs will provide examples in which several structural properties occur together.

Where possible, comparisons will hold node count, edge count, and feature dimensions constant while varying connectivity. This will help distinguish effects associated with topology from those caused simply by increasing the amount of data.

Profiling will determine which operations account for substantial execution time and how their behavior changes across graphs. Relevant observations may include uneven work distribution, repeated feature accesses, or contention during aggregation. These observations will guide the choice of optimization target.

\subsection{Developing and Applying Specialization}

After identifying a measurable bottleneck, the project will implement a focused alternative for a supported operation. Candidate approaches may include a different sparse layout, a different assignment of work to threads, or combining operations to reduce intermediate memory traffic.

Specialized GPU implementations will use Triton where appropriate. CUDA/C++ may be considered if the selected operation requires control that is difficult to express effectively in Triton.

Each candidate will first be tested independently, then evaluated within a complete inference workload. Graph analysis and preprocessing will be measured separately from repeated execution.

The project will test whether graph characteristics predict when a candidate is beneficial, potentially using a simple selection rule. It will then examine how to connect successful implementations to the \texttt{torch.compile} workflow, while retaining TorchInductor for other operations.
